% !TeX program = lualatex
% Compile twice: lualatex 133.tex
% All references and prose are embedded; no BibTeX database or external inputs.
\documentclass[11pt,a4paper]{article}
\usepackage[margin=24mm,headheight=14pt]{geometry}
\usepackage{fontspec}
\setmainfont{Latin Modern Roman}
\setsansfont{Latin Modern Sans}
\setmonofont{Latin Modern Mono}
\usepackage{amsmath,amssymb,mathtools}
\usepackage{unicode-math}
\setmathfont{Latin Modern Math}
\usepackage{microtype}
\usepackage{enumitem,xurl,needspace}
\usepackage[dvipsnames]{xcolor}
\usepackage{hyperref}
\hypersetup{unicode=true,colorlinks=true,linkcolor=MidnightBlue,urlcolor=MidnightBlue,
 pdftitle={500 Mathematical Problem Statements},pdfauthor={Source-annotated compilation},bookmarksnumbered=true}
\usepackage{bookmark}
\usepackage{tocloft}
\setlength{\cftsecnumwidth}{3em}
\cftsetpnumwidth{2.5em}
\cftsetrmarg{4em}
\usepackage{titlesec}
\titleformat{\section}{\Large\bfseries\raggedright}{\thesection}{0.7em}{}
\usepackage{fancyhdr}
\pagestyle{fancy}\fancyhf{}
\fancyhead[L]{\small\sffamily Mathematical problem statements}
\fancyhead[R]{\small Source edition 22}
\fancyfoot[C]{\thepage}
\setlength{\parindent}{0pt}
\setlength{\parskip}{5pt plus 1pt}
\setlength{\emergencystretch}{3em}
\setcounter{tocdepth}{1}
\setcounter{secnumdepth}{1}
\setlist[itemize]{leftmargin=3em,itemsep=5pt,topsep=4pt,parsep=0pt}
\allowdisplaybreaks
\newcommand{\N}{\mathbb{N}}
\newcommand{\Z}{\mathbb{Z}}
\newcommand{\Q}{\mathbb{Q}}
\newcommand{\R}{\mathbb{R}}
\newcommand{\C}{\mathbb{C}}
\newcommand{\F}{\mathbb{F}}
\newcommand{\A}{\mathbb{A}}
\newcommand{\PP}{\mathbb P}
\newcommand{\E}{\mathbb{E}}
\newcommand{\eps}{\varepsilon}
\newcommand{\dd}{\,\mathrm d}
\newcommand{\sref}[2]{\hyperlink{source:#1:#2}{[S#2]}}
\newcommand{\eref}[2]{\hyperlink{extra:#1:#2}{[E#2]}}
% Turn the next switch off for a shorter reading copy. The quotations preserve
% every exactTarget field, including qualifications omitted from a short summary.
\newif\ifshowcatalogue
\showcataloguetrue
\newcommand{\cataloguescope}[1]{\ifshowcatalogue
 \paragraph{Exact scope in the supplied catalogue.}
 \begin{quote}\small #1\end{quote}\fi}
\hypersetup{pdftitle={Problem 133 - Normality of algebraic irrationals},pdfauthor={Open Problems Math research notebook}}
\fancyhead[L]{\small\sffamily Open Problems Math \textbar\ Research notebook}
\fancyhead[R]{\small\sffamily 133 \textbar\ 27 September 2026}
\numberwithin{equation}{section}
\begin{document}
\setcounter{section}{132}
\setlength{\parskip}{3pt plus 1pt}
\setlist[itemize]{itemsep=3pt,topsep=3pt}
\titlespacing*{\paragraph}{0pt}{1.5ex plus .5ex minus .2ex}{1em}
% ================================================================
% problemId: problem.borel-normal-number-conjecture-for-irrational-algebraic-numbers
% source releaseRank: 133; source releaseStatus: open
% exactTarget SHA-256: 4e3102b938353c1f4f811471a84f80e55f53c7b8deb785b753ae907159bbf17b
\clearpage
\section{\texorpdfstring{Borel normal number conjecture for irrational algebraic numbers}{Borel normal number conjecture for irrational algebraic numbers}}\label{133}
\subsection{Definitions and mathematical statement}
Write the fractional part of a real number in base $b\ge2$ as $0.d_1d_2\ldots$. It is normal to base $b$ if for every length $k$ and every word $w\in\{0,\ldots,b-1\}^k$,
\[
 \lim_{N\to\infty}\frac{\#\{1\le j\le N:(d_j,\ldots,d_{j+k-1})=w\}}{N}=b^{-k}.
\]
The conjecture requires this for every real algebraic irrational number and every integer base. Algebraic means a root of a nonzero rational polynomial. Frequency of individual digits alone is weaker than normality of all blocks.
\par\smallskip\textit{Formulation and concept references:} \sref{133}{1}.
\cataloguescope{Prove or disprove that every real irrational algebraic number is normal to every integer base b >= 2.}
\subsection{Short English statement}
Do the digits of every algebraic irrational have the uniform block frequencies of a normal number in every integer base?
% BEGIN RESEARCH ADDITION 133
\subsection{Research attempt}

\paragraph{Current frontier and scope.}
The target is normality of every real algebraic irrational in every
integer base, including all finite digit blocks.  The source discusses
the distance between known transcendence and digit-structure theorems
and this frequency assertion. \sref{133}{1}
A relevant established result is Adamczewski--Bugeaud's theorem that
the block complexity $p(\ell)$ of an algebraic irrational's base-$b$
expansion satisfies $p(\ell)/\ell\to\infty$.  This counts \emph{distinct}
blocks and is not a statement about their frequencies.
\eref{133}{1}  The search through 27 September 2026 did not provide a
verified theorem proving the full normality statement.  The older
structural theorem is used below as a benchmark, not represented as
a complete current inventory of digit estimates.

\paragraph{Attack map.}
A Diophantine-approximation route tries to turn a biased digit expansion
into excessively good rational approximations.  Long constant runs
can indeed be bounded, but frequency bias need not create them.
An exponential-sum route uses algebraicity in a cancellation estimate
for a lacunary sequence; after differencing, the same difficult type
of sum reappears.  A combinatorial route would strengthen the known
block-complexity theorem into uniform frequencies, but maximal block
complexity by itself is insufficient.  These routes are compared using
explicit deductions and counterexamples to their proposed intermediate
inferences.

\paragraph{Attempt I: the exact exponential-sum target.}
Replace $\alpha$ by its fractional part, which remains an algebraic
irrational in $(0,1)$.  Fix a base $b\geq2$.  For
$x_n=\{b^n\alpha\}$, a word of length $\ell$ representing the integer
$j\in\{0,\ldots,b^\ell-1\}$ starts at digit $n+1$ exactly when
\[
 x_n\in[j/b^\ell,(j+1)/b^\ell).
\]
No $x_n$ is an endpoint, since that would make $\alpha$ rational.
Consequently equidistribution of $(x_n)$ implies normality.  Conversely,
if all these block frequencies are correct, approximate any interval
from inside and outside by unions of base-$b$ intervals of mesh
$b^{-\ell}$.  Their total length difference tends to zero, giving
uniform distribution in every interval.  Thus the two assertions are
equivalent, with no density-one subsequence substituted.

The corresponding Fourier criterion is
\begin{equation}\label{133:weyl}
 \forall h\in\Z\setminus\{0\}:\qquad
 S_h(N):=\sum_{n=0}^{N-1}\exp(2\pi i h b^n\alpha)=o(N).
\end{equation}
To recall why, equidistribution gives the averages of each nonconstant
Fourier character equal to its integral zero.  Conversely, Fourier
averages determine averages of trigonometric polynomials; uniform
approximation of continuous periodic functions by such polynomials,
followed by continuous upper and lower approximations to interval
indicators, gives equidistribution.  These are exact logical
reductions, not a new cancellation theorem.  The source's discussion
of uniform digit distribution provides the original context.
\sref{133}{1}

\paragraph{Attempt II: differencing and its circular stopping point.}
Put $z_n=\exp(2\pi i h b^n\alpha)$ for $0\leq n<N$ and extend it by
zero outside that range.  For $1\leq H\leq N$, summing overlapping
blocks and applying Cauchy--Schwarz gives
\begin{equation}\label{133:vdc}
 |S_h(N)|^2\leq\frac{N+H-1}{H^2}
 \left(HN+2\sum_{d=1}^{H-1}(H-d)\Re C_d(N)\right),
\end{equation}
where
\[
 C_d(N)=\sum_{n=0}^{N-d-1}z_{n+d}\overline{z_n}
       =\sum_{n=0}^{N-d-1}
           \exp(2\pi i h(b^d-1)b^n\alpha).
\]
For the inequality, use
$HS_h(N)=\sum_{r=-(H-1)}^{N-1}\sum_{j=0}^{H-1}z_{r+j}$.
Expanding the squared inner sums counts a pair separated by $d$
exactly $H-d$ times, which proves \eqref{133:vdc} including its
endpoints and constant factors.

For any fixed $H$, bounds $C_d(N)=o(N)$ for $1\leq d<H$ would imply
$\limsup_N|S_h(N)|^2/N^2\leq1/H$; then $H\to\infty$ would prove
\eqref{133:weyl}.  But $C_d(N)$ is precisely another sum of the same
form, with frequency $h(b^d-1)$.  Multiplying an algebraic irrational
by this nonzero integer preserves algebraicity and irrationality;
it supplies no quantitative cancellation.  Invoking normality of
that new algebraic number to bound the correlation would assume an
instance of the very universal conjecture being attempted.  The
useful missing lemma is an independent bound for these correlations,
not the formal differencing identity.

\paragraph{Attempt III: what Roth-type approximation actually forbids.}
Suppose that after its first $m$ base-$b$ digits, $\alpha$ has $L$
consecutive zeros.  If the prefix represents $p/b^m$, then
\[
 0<|\alpha-p/b^m|\leq b^{-m-L}.
\]
For every fixed $\varepsilon>0$, Roth's approximation theorem implies
a constant $c_{\alpha,\varepsilon}>0$ such that every reduced rational
$p'/q$ satisfies
$|\alpha-p'/q|\geq c_{\alpha,\varepsilon}q^{-2-\varepsilon}$.
The finitely many exceptional approximants are incorporated by reducing
the positive constant.  This theorem and its digit applications are
part of the source's Diophantine-approximation discussion.
\sref{133}{1}

Reducing $p/b^m$ gives $q\leq b^m$, so the lower bound is at least
$c_{\alpha,\varepsilon}b^{-(2+\varepsilon)m}$.  Combining the two
inequalities yields
\begin{equation}\label{133:runs}
 L\leq(1+\varepsilon)m+O_{\alpha,b,\varepsilon}(1).
\end{equation}
A run of digits all equal to $b-1$ is treated with the upper rational
endpoint $(p+1)/b^m$.  Its difference from $\alpha$ is nonzero for
an irrational number, and the same bound follows.  No effective value
for Roth's constant is asserted.

This proof rules out certain very long uniform runs.  It does not
turn a fixed positive digit-frequency bias into a contradiction:
a biased sequence can change digits often enough to have no long
runs at all.  Nor does the block-complexity theorem
$p(\ell)/\ell\to\infty$ provide an estimate for the relative
frequencies required by normality. \eref{133}{1}

\paragraph{A probabilistic stress test of both structural shortcuts.}
Consider independent binary digits with probabilities
$\Pr[d_n=1]=1/3$ and $\Pr[d_n=0]=2/3$.  With probability one, the
frequency of ones is $1/3$, by the strong law, so the resulting number
is not normal to base two.  Every fixed finite word has positive
probability at each of infinitely many disjoint blocks.  Independence
shows that it appears infinitely often almost surely.  Intersecting
these probability-one events over the countably many words gives
\[
 p(\ell)=2^\ell\quad\text{for every }\ell
\]
almost surely: the block complexity is maximal, not merely superlinear.

The same sequences almost surely have only logarithmically long
constant runs.  Among the first $N$ digits, the probability of a
constant run of length $L$ is at most
$2N(2/3)^L$.  Take $N=2^j$ and $L=\lceil C\log N\rceil$ with
$C$ sufficiently large; the resulting probabilities sum to a finite
number over $j$.  The union bound and the elementary first
Borel--Cantelli lemma then give a logarithmic bound for all large
dyadic $N$, and hence for all large $N$.  This is far stronger than
\eqref{133:runs} while still compatible with persistent frequency bias.

These numbers are not proposed algebraic counterexamples.  The
construction demonstrates that the two \emph{consequences} extracted
above---many distinct blocks and no excessively long constant
runs---do not logically imply normality, even together.  A successful
proof must exploit further arithmetic information about algebraicity,
not just these combinatorial consequences.  Conversely, no theorem
here says that every biased sequence satisfying these conditions is
transcendental.

\paragraph{Exact finite digit experiments.}
To check digit-frequency computations without floating-point ambiguity,
let $N\geq1$ and use the integer
\[
 M=\left\lfloor\sqrt{2b^{2N}}\right\rfloor.
\]
Its defining certificate is the exact pair of inequalities
$M^2\leq2b^{2N}<(M+1)^2$.  The $N$ base-$b$ digits of $M-b^N$,
padded with leading zeros, are exactly the first $N$ fractional digits
of $\sqrt2$.  The companion script applies this to a binary prefix
of length $4096$ and a decimal prefix of length $1000$, and records
all overlapping block counts for lengths one through three.

For these finite windows there are $N-\ell+1$ complete length-$\ell$
blocks, not $N$ blocks with a silently omitted suffix.  The report
states that denominator explicitly.  The integer-square certificate
proves that the tested digits are correct; it says nothing about the
asymptotic behavior of their frequencies, or about another algebraic
number or another base.

The binary prefix has $2066$ zeros and $2030$ ones. Its $4094$
complete overlapping triples have the following exact counts:
\[
\begin{array}{c|r@{\qquad}c|r}
\text{word}&\text{count}&\text{word}&\text{count}\\\hline
000&555&100&516\\
001&516&101&477\\
010&501&110&493\\
011&493&111&543
\end{array}
\]
They sum to $4094$, not $4096$; the largest absolute deviation from
frequency $1/8$ is exactly $173/16376$. The decimal prefix yields
\[
\begin{array}{c|rrrrrrrrrr}
\text{digit}&0&1&2&3&4&5&6&7&8&9\\\hline
\text{count}&108&98&109&82&100&104&90&104&113&92
\end{array}
\]
with denominator $1000$. These outputs illustrate the difference
between an exact finite-prefix certificate and the all-length limiting
assertion: even correct counts for all words in this finite window do
not bound the discrepancy at later positions. The full two- and
three-digit decimal tables are retained in the machine-readable report.

\paragraph{Outcome.}
\textbf{Outcome: Exploratory approach with unresolved gap.}

The Fourier equivalence, differencing identity, and long-run bound are
proved.  The probabilistic example refutes the proposed implication
from maximal block complexity and short constant runs to normality.
The arithmetic cancellation estimate \eqref{133:weyl} remains unproved;
no algebraic counterexample, resolution, or historical novelty is claimed.
% END RESEARCH ADDITION 133
\subsection{Sources}
\begin{itemize}\raggedright
\item[\textnormal{[S1]}]\hypertarget{source:133:1}{} Michel Waldschmidt, Words and Transcendence, 2005.
\url{https://webusers.imj-prg.fr/~michel.waldschmidt/articles/pdf/WordsTranscendence.pdf}.
{\small\textit{Catalogue formulation source.}}
% BEGIN RESEARCH REFERENCES 133
\item[\textnormal{[E1]}]\hypertarget{extra:133:1}{} Boris Adamczewski and Yann Bugeaud.
\emph{On the complexity of algebraic numbers I. Expansions in integer bases}, Annals of Mathematics \textbf{165} (2007), no.~2, 547--565, Theorem 1 (block complexity, not frequency normality).
\url{https://arxiv.org/html/math/0511674}.
{\small\textit{Consulted 27 September 2026 for the indicated result or scope.}}
% END RESEARCH REFERENCES 133
\end{itemize}


\end{document}
